\section{Preliminaries}

% \dotparagraph{Notation}
\textbf{Notation:} We denote by \BN the set of natural numbers including $0$.
Letting $L$ be a bounded sequence, we denote by $L[i]$ the $i^\text{th}$ element (starting from $0$),
and by $|L|$ its length.
Negative indices address elements from the end, so $L[-i]$ is the $i^\text{th}$ element from
the end, and $L[-1]$ in particular is the last.
%We use $L[i{:}j]$ to denote the subarray of $L$
% consisting of the elements indexed from $i$ (inclusive) to $j$ (exclusive).
The notation $L[i{:}]$ means the
subarray of $L$ from $i$ onwards, while $L[{:}j]$ means the subsequence of $L$ up to (but not including) $j$.
We denote an empty sequence by $[\,]$.
We denote the concatenation of sequences $L_1$ and $L_2$ by $L_1 \concat L_2$.

\textbf{Pseudocode notation:} Command `\textbf{require} $P$' causes a function to terminate immediately and return $\false$ if $P$ evaluates to $\false$.
Notation `\textbf{upon} $e$' causes a block of code to be executed when event $e$ occurs.
Notations `$\msg{MSG}{} \leftarrow p$' and `$\msg{MSG}{} \rightarrow p$' denote receiving and sending a message $MSG$ from and to party $p$, respectively.
Finally, $x: a \in A \rightarrow b \in B$ denotes that variable $x$ is a map from elements of type $A$ to elements of type $B$. When obvious from the context, we do not explicitly write the types $A$ or $B$.
For a map $x$, the operations $x.\opKeys$ and $x.\opValues$ return all keys and all values in $x$, respectively.
With $\emptyset$ we denote an empty map.
% \subsection{Execution Model}
\label{sec:system}

\textbf{Parties:}
We consider $n$ \emph{replicas} $\replicas=\{\replicas_1,\ldots,\replicas_n\}$ and an unknown number of \emph{clients}. Parties are \emph{stateful}, \textit{i.e.}, store \emph{state} between executions of different algorithms.
We assume that replicas are known to all parties and register their public keys (for which they have corresponding secret keys) in a Public Key Infrastructure (PKI). Clients do not register keys in the PKI.
% \vspace{-3mm}

\textbf{Adversarial model:}
We call a party (replica or client) \emph{honest}, if it follows the protocol, and \emph{malicious} otherwise.
% Our constructions assume limits on \alpha and \beta, but make no honesty assumptions on clients.
We assume \emph{static corruptions},  \textit{i.e.}, the set of malicious replicas is decided before the execution starts and remains constant.
This work uses a combination of two adversarial models, the \emph{Byzantine} and the \emph{omission} models.
In the \emph{Byzantine} model, corrupted replicas are malicious and may deviate arbitrarily from the protocol. The adversary has access to the internal state and secret keys of all corrupted parties. We denote by $\beta \in \BN$ the  number of Byzantine replicas in an execution.
The Byzantine adversary is modelled as a probabilistic polynomial time overarching entity that is invoked in the stead of every corrupted party.
That is, whenever the turn of a corrupted party comes to be invoked by the environment, the adversary is invoked instead.
In the \emph{omission} model, corrupted replicas may only deviate from the protocol by dropping messages that they were supposed to send, but follow the protocol otherwise. Observe that this includes crash faults, where replicas crash (\textit{i.e.} stop execution) and remain crashed until the end of the execution of an algorithm. We denote by $\gamma \in \BN$ the number of omission-faulty replicas in an execution.

% \oa{We allow \emph{dynamic corruptions}, i.e., the set of malicious replicas can differ in each round, as long as not larger than \beta.}

%\vspace{-3mm}
%\dotparagraph{Stateful clients} A stateful client is a party who stores \emph{state} (\textit{i.e.}, data) between executions of different algorithms.
% \vspace{-3mm}

\textbf{Modeling time:}
We assume that time proceeds in discrete \emph{rounds}, and parties have clocks allowing them to determine the current round.
For simplicity, our analysis will assume \emph{synchronized clocks}.
Notice that although we assume synchronized clocks as a setup, clock synchronization can be achieved in partially synchronous networks~\cite{psync} using existing techniques~\cite{clockoverview}, also in the case where replicas gradually join the network~\cite{bootclocksync}.
% We denote by $\tzero$ the round at which the execution starts.
By \emph{timestamp} we refer to a round number assigned to some event.

\textbf{Modeling network:}
We denote by $\delta \in \CN$ the actual delay (measured in number of rounds) it takes to deliver a message between two honest parties,
a number which is \emph{finite} but \emph{unknown} to all parties.
We denote by $\Delta \in \CN$ an upper bound on this delay, \textit{i.e.}, $\delta \leq \Delta$, which is also \emph{finite}.
In the \emph{synchronous} model, $\Delta$ is \emph{known} to all parties.
% (namely, the honest code is parametrized by $\Delta$).
In the \emph{partially synchronous} model~\cite{psync}, $\Delta$ is \emph{unknown} but still finite, \textit{i.e.}, all messages are eventually delivered.
%The environment is constrained to deliver all honest network messages within $\delta$. Namely, for each message $m$ sent
%by an honest party at round $r$, the environment observes and remembers the inboxes (decided by the adversary)
%of all honest parties during rounds $r + 1, \ldots, r + \delta$. If none of these inboxes contains $m$, the environment
%appends $m$ to the network inbox of the recieving party prior to invoking it at round $r + \delta$, thereby forcing the
%inclusion of this message.
% In this work, we present a general protocol, which we then analyze in the \emph{synchronous} and in the \emph{asynchronous} models, providing different guarantees in each world.
% We now define responsive protocols, which is the case of our main protocol \podtm-core.
A protocol is called \emph{responsive} if it does not rely on knowledge of $\Delta$ and its liveness guarantees depend only on the actual network delay $\delta$.
%this is very fluffy, we should make it precise for the DISC submission.

% \subsection{Cryptographic primitives}
\dotparagraph{Digital signatures} We assume that replicas (and auctioneers in \bidsettm-core) authenticate their messages with digital signatures. A digital signature scheme consists of the following three algorithms, satisfying the EUF-CMA security~\cite{GMR88}:
(1) $\textsl{KeyGen}(1^{\kappa})$:
The key generation algorithm takes as input a security parameter $\kappa$ and outputs a secret key $\sk$ and a public key $\pk$.
(2) $\opSign{\sk, m} \rightarrow \sig$:
The signing algorithm takes as input a private key
\sk and a message $m \in \zo^*$ and returns a signature $\sig$.
(3)$\opVerify{\pk, m, \sig} \rightarrow b \in \zo$:
The verification algorithm takes as input a public
key $\pk$, a message $m$, and a signature $\sig$, and outputs a bit $b\in \zo$.
We say \sig is a \emph{valid} signature on $m$ with respect to \pk if $\opVerify{pk, m, \sig} = 1$.





\dotparagraph{Accountable safety}
\label{sec:accountable-safety}
Taking a similar approach as Neu, Tas, and Tse~\cite[Def. 4]{FC:NeuTasTse22}, we define \emph{accountable safety} through an \emph{identification function}.

\begin{definition}[Transcript and partial transcript]
    We define as \emph{transcript} the set of all network messages sent by all parties in an execution of a protocol.
    A \emph{partial transcript} is a subset of a transcript.
\end{definition}

\begin{definition}[\beta-Accountable safety]\label{def:accountable-safety}
    A protocol satisfies \emph{accountable safety} with \emph{resilience} \beta if its interface contains a function $\identify{\tran} \rightarrow \cheaters$,
    which takes as input a partial transcript $\tran$ and outputs a set of replicas $\cheaters \subset \replicas$, such that the
    following conditions hold except with negligible probability.
    \begin{description}
        \item[Correctness:] If safety is violated, then there exists a partial transcript \tran, such that $\identify{\tran} \rightarrow \cheaters$ and $|\cheaters| > \beta$.
        \item[No-framing:] For any partial transcript \tran produced during an execution of the protocol, the output of $\identify{\tran}$ does not contain honest replicas.
    \end{description}
\end{definition}

\begin{remark}
    For simplicity, we have defined the transcript based on messages sent by \emph{all} replicas.
    We can also define a \emph{local transcript} as the set of messages observed by a single party.
    As will become evident from the implementation of \identify{}, in practice, adversarial behavior can be identified from the local transcripts of a single party or  of a pair of parties.
    %This optimization can be used in the implementation of a protocol.
\end{remark}